% !PRINTED 23
\weberpage{23}
%<<CONT
and since $R$ is now a greatest normal divisor of $Q$ and of $Q'$,
we can form from this two new composition series of $P$,
namely
\weq{26}{\begin{aligned}
&P,\ Q,\ R,\ R_1,\ R_2,\ R_3 \ldots \\
&\hphantom{P,\ }\nu,\ \nu',\ \mu_1,\ \mu_2,\ \mu_3 \ldots
\end{aligned}}
\weq{27}{\begin{aligned}
&P,\ Q',\ R,\ R_1,\ R_2,\ R_3 \ldots \\
&\hphantom{P,\ }\nu',\ \nu,\ \mu_1,\ \mu_2,\ \mu_3 \ldots
\end{aligned}}
The two composition series (26) and (27) of $P$ therefore have
the same index series.

If we now assume that the theorem to be proved has already been
established as correct for groups whose degree is $\leqq \frac{1}{2}n$
(where $n$ denotes the degree of $P$), it follows that all index
series of $Q$, whose degree is after all a proper divisor of $n$
and hence at most $= \frac{1}{2}n$, agree with one another, and hence
that the series
\[
\nu',\ \mu_1,\ \mu_2,\ \mu_3 \ldots
\]
\[
\nu_1,\ \nu_2,\ \nu_3,\ \nu_4 \ldots,
\]
agree, apart from the arrangement. In the same way the index
series of $Q'$
\[
\nu,\ \mu_1,\ \mu_2,\ \mu_3 \ldots
\]
\[
\nu_1',\ \nu_2',\ \nu_3',\ \nu_4' \ldots
\]
agree. Hence the two index series of $P$
\[
\nu,\ \nu_1,\ \nu_2,\ \nu_3,\ \nu_4 \ldots
\]
\[
\nu',\ \nu_1',\ \nu_2',\ \nu_3',\ \nu_4' \ldots
\]
also agree with one another, since the first agrees with $\nu, \nu', \mu_1, \mu_2, \mu_3
\ldots$, and the second with $\nu', \nu, \mu_1, \mu_2, \mu_3 \ldots$

For groups of degree 2, however, which have only the single normal
divisor 1, the theorem is evident, and it is therefore proved in
general by complete induction.

If in two composition series of $P$, which together with their
index series we shall now represent thus
\weq{28}{\begin{aligned}
&P,\ P_1,\ P_2,\ P_3 \ldots P_{\mu-1},\ 1 \\
&\hphantom{P,\ }j_1,\ j_2,\ j_3 \ldots j_{\mu-1},\ j_\mu
\end{aligned}}
\weq{29}{\begin{aligned}
&P,\ P_1',\ P_2',\ P_3' \ldots P_{\mu-1}',\ 1 \\
&\hphantom{P,\ }j_1',\ j_2',\ j_3' \ldots j_{\mu-1}',\ j_\mu'
\end{aligned}}
a common term $P_r = P_s'$ occurs, then the theorem on the
constancy of the index series holds also for the group $P_r =
P_s'$, and from this it follows first that we must have $s = r$, and
%CONT>>
